%==============================================================================
%  Valter Nunez — CV en español (México)  ·  1 página, compatible con ATS
%  Compilar en Overleaf (XeLaTeX) o localmente: `tectonic Valter-Nunez-CV-ES.tex`
%==============================================================================
\documentclass[10pt,a4paper]{article}

\usepackage[T1]{fontenc}
\usepackage[utf8]{inputenc}
\usepackage[default]{lato}
\renewcommand*\familydefault{\sfdefault}
\usepackage[margin=1.05cm]{geometry}
\usepackage{enumitem}
\usepackage{titlesec}
\usepackage{xcolor}
\usepackage[hidelinks]{hyperref}

\definecolor{ink}{HTML}{0A0A0A}
\definecolor{accent}{HTML}{2563EB}
\definecolor{soft}{HTML}{45474B}
\color{ink}
\hypersetup{colorlinks=true, urlcolor=accent}

\setlength{\parindent}{0pt}
\linespread{1.04}
\pagestyle{empty}

\titleformat{\section}
  {\large\bfseries\color{ink}}{}{0em}{}
  [\vspace{-0.55em}\color{accent}\rule{\linewidth}{1.1pt}]
\titlespacing*{\section}{0pt}{0.5em}{0.4em}

\newcommand{\rol}[4]{%
  \textbf{#1} \hfill \textbf{#3}\\[1pt]
  {\color{soft}\textit{#2} \hfill \textit{#4}}\par\vspace{2pt}}
\newcommand{\proyecto}[2]{\textbf{#1} \,---\, #2\par\vspace{1pt}}
\newcommand{\hab}[2]{\textbf{\color{soft}#1:}\ #2\par\vspace{1.5pt}}

\newlist{tight}{itemize}{1}
\setlist[tight]{leftmargin=1.1em, itemsep=1.6pt, topsep=2pt, parsep=0pt,
  label=\textcolor{accent}{\small$\bullet$}}

\begin{document}

%--- ENCABEZADO ---------------------------------------------------------------
{\Huge\bfseries Valter Nunez}\\[3pt]
{\color{soft}\large Ingeniero Backend y de Sistemas Distribuidos}\\[5pt]
{\small
\href{mailto:valter.nunez@outlook.com}{valter.nunez@outlook.com} \quad\textcolor{soft}{|}\quad
\href{https://www.linkedin.com/in/valternunez}{linkedin.com/in/valternunez} \quad\textcolor{soft}{|}\quad
\href{https://valternunez.github.io}{valternunez.github.io}
}

%--- RESUMEN ------------------------------------------------------------------
\section*{Resumen}
Ingeniero backend con cinco años en e-commerce de alto tráfico: sistemas orientados a
eventos y tolerantes a fallos. Desacoplé el \textit{83\%} de los flujos de venta críticos del ERP y
reduje el gasto total de envíos un \textit{\textasciitilde78\%} con enrutamiento automático al menor
costo. Actualmente curso la \textbf{Maestría en Ciencia de Datos en LSE}, en Londres. Elegible para
la Graduate visa del Reino Unido (18 meses, sin necesidad de patrocinio); en busca de patrocinio Skilled Worker para permanecer a largo plazo.

%--- EXPERIENCIA --------------------------------------------------------------
\section*{Experiencia}

\rol{Ingeniero Backend}{ZeBrands — e-commerce de alto tráfico}{Sep 2021 -- Mar 2026}{Remoto, México}
\begin{tight}
  \item Diseñé y construí \textbf{OZ}, una plataforma orientada a eventos que sostiene las ventas durante caídas del ERP, desacoplando el \textit{83\%} de los flujos de venta críticos de llamadas síncronas al ERP; el análisis de modos de falla proyectó una reducción del \textit{91\%} en el tiempo de caída que bloquea ventas (FastAPI, PostgreSQL, RabbitMQ, Redis); hoy en producción en \textit{\textasciitilde195 tiendas} (México, en expansión a Brasil), sosteniendo \textit{\textasciitilde MX\$102M} en ventas durante el pico de nueve días del Hot Sale.
  \item Construí un middleware de envíos que unifica las APIs de varias paqueterías (FedEx, DHL) con enrutamiento automático al menor costo, reduciendo el gasto total de envíos un \textit{\textasciitilde78\%} ($\approx$ \$500K USD/año).
  \item Migré la analítica de reseñas a un microservicio dedicado, mejorando la disponibilidad de datos y el desempeño de las consultas.
  \item Mantuve microservicios existentes de cara al cliente, módulos del ERP y APIs GraphQL (Frappe, TypeScript), y monté logging/observabilidad en Grafana que redujo los tiempos de detección y respuesta.
  \item Mentoré a un desarrollador backend junior; contribuí a las revisiones de arquitectura y a los estándares de código backend de los equipos de cliente y promociones.
\end{tight}
\vspace{2pt}

\rol{Cofundador y Desarrollador}{Laiter — estudio de software}{2017 -- 2019}{México}
\begin{tight}
  \item Entregué aplicaciones web para clientes de principio a fin (Django, JavaScript) y construí los primeros tableros de analítica de uso con los que los clientes tomaban decisiones; me hice cargo de los despliegues y del soporte en producción.
\end{tight}

%--- PROYECTOS ----------------------------------------------------------------
\section*{Proyectos seleccionados}
\begin{tight}
  \item \proyecto{Nexus}{arquitecto único de un SaaS multiinquilino para administración de inmuebles: motor financiero de partida doble, Stripe Connect (cargos, SPEI, pagos, disputas), RBAC por inquilino, un asistente con LLM (Gemini) con medición de uso, límites de tasa y salida estructurada, más un panel de administración en Next.js y una PWA para residentes. Construido en solitario mediante un flujo de desarrollo multiagente orquestado (Claude Code, git worktrees). En desarrollo activo (FastAPI, async SQLAlchemy).}
  \item \proyecto{Estudio de atención del Balón de Oro}{descomposición bayesiana (PyMC, \textasciitilde210 temporadas-jugador) que aísla un residual inexplicado de atención mediática tras controlar por mérito y éxito del equipo; sitio de scrollytelling en D3 + reporte en Quarto.}
  \item \proyecto{Momentum del Mundial 2026}{pipeline de datos de experimento natural (scraping local $\rightarrow$ GitHub Actions $\rightarrow$ Pages) que se regenera a diario.}
\end{tight}

%--- EDUCACIÓN ----------------------------------------------------------------
\section*{Educación}
\rol{Maestría en Ciencia de Datos}{London School of Economics (LSE)}{2026 -- 2027}{En curso, Londres}
\vspace{2pt}
\rol{Ing. en Ciencias y Tecnologías Computacionales}{Tecnológico de Monterrey (ITESM)}{2016 -- 2021}{Prom. 90/100}
\vspace{1pt}
{\color{soft}\small Semestre de intercambio en RMIT University, Melbourne (Prom. 94/100). Certificado Profesional en Ciencia de Datos de IBM (en curso).}

%--- HABILIDADES --------------------------------------------------------------
\section*{Habilidades técnicas}
\hab{Lenguajes}{Python, TypeScript, JavaScript, SQL \;(familiar: Go)}
\hab{Backend}{FastAPI, Django, GraphQL, Celery, Frappe/ERPNext}
\hab{Datos e infraestructura}{PostgreSQL, Redis, RabbitMQ, MongoDB, Docker, Grafana, GitHub Actions}
\hab{Ciencia de datos}{pandas, NumPy, scikit-learn, PyMC, inferencia bayesiana, Quarto}
\hab{IA / LLM}{Integración de LLM (Gemini), medición de uso y límites de tasa, salida estructurada, desarrollo orquestado con agentes}
\hab{Prácticas}{Diseño orientado a eventos y de microservicios, pruebas, observabilidad, revisión de código}
\hab{Idiomas}{Español (nativo), Inglés (C2), Alemán (A2), Italiano (A1)}

\end{document}
